% --- Archon physics formalization source begin ---
% archon:physics
% archon:covers PhyXMiniProblems/problem_phyx_mini_0100.lean
% archon:source-report reports/phyx_mini/problem_phyx_mini_0100.source.json

\chapter{Physics problem phyx\_mini\_0100}
\label{ch:PhyXMiniProblems_problem_phyx_mini_0100}

\paragraph{Problem source.}
\#\# Physical scenario

Light is incident normally on the short face of a \$30\textasciicircum{}\textbackslash{}circ\$-\$60\textasciicircum{}\textbackslash{}circ\$-\$90\textasciicircum{}\textbackslash{}circ\$ prism (figure). A drop of liquid is placed on the hypotenuse of the prism.

\#\# Question

If the index of refraction of the prism is 1.56, find the maximum index that the liquid may have for the light to be totally reflected.

\#\# Answer choices

- A: A: 1.35

- B: B: 2.43

- C: C: 0.94

- D: D: 1.19

\#\# Figure caption (auxiliary; use the image as primary evidence)

The image depicts a right triangle with angles marked as 90°, 60°, and 30°. The 90° angle is indicated by a small square in the corner of the triangle. 

There is an arc above the 60° angle, highlighting it. A curved arrow is associated with the 30° angle, pointing upwards. 

Inside the triangle, there is a purple arrow starting from the left side, proceeding diagonally, and splitting into two directions: one parallel to the hypotenuse and another moving towards the 30° angle. The arrow emphasizes a path or direction within the triangle.

\#\# Dataset metadata

Geometrical Optics; Physical Model Grounding Reasoning, Spatial Relation Reasoning

\paragraph{Recorded answer/context.}
A: A: 1.35

\paragraph{Figure/image path.}
/root/proposal\_for\_physic/science-mango/phyx\_mini\_run/phyx\_data/test\_image/100.png

\paragraph{Formalization target.}
create a compiling Lean file with sorry bodies at `PhyXMiniProblems/problem\_phyx\_mini\_0100.lean`. The Lean declarations must preserve the physical quantities, dimensions or dimensional roles, figure labels, governing-law hypotheses, and final relation expressed by this problem.
Use Mathlib/Physlib names found through LeanExplore where available. If a physics API is missing, introduce faithful local abstractions rather than scalar placeholder aliases.

\begin{theorem}[Physics formalization target]
\label{thm:physics:phyx_mini_0100:target}
\lean{PhyXMiniProblems.ProblemPhyXMini0100.maximum_liquid_refractive_index}
\uses{def:physics:phyx-mini-0100:phyxminiproblems-problemphyxmini0100-prismliquidsetup, def:physics:phyx-mini-0100:phyxminiproblems-problemphyxmini0100-matchesprismfigure, def:physics:phyx-mini-0100:phyxminiproblems-problemphyxmini0100-matchesprismindexreadout, def:physics:phyx-mini-0100:phyxminiproblems-problemphyxmini0100-satisfiesprismraygeometry, def:physics:phyx-mini-0100:phyxminiproblems-problemphyxmini0100-hasphysicalprismparameters, def:physics:phyx-mini-0100:phyxminiproblems-problemphyxmini0100-admissibleliquidrefractiveindices, def:physics:phyx-mini-0100:phyxminiproblems-problemphyxmini0100-criticalliquidindexcandidate, def:physics:phyx-mini-0100:phyxminiproblems-problemphyxmini0100-isclosestanswerchoice}
The assigned autoformalize agent should translate this physics problem into Lean declarations in the covered file, with theorem and lemma proof bodies written as `by sorry`.
\end{theorem}
\begin{proof}
This is an autoformalization task, not a proof task. Produce statements that can later be proved by the physics prover without weakening the physical model.
\end{proof}

% --- Archon declaration topology begin ---
\section{Declaration topology}
\begin{definition}[Prism Vertex]
  \label{def:physics:phyx-mini-0100:phyxminiproblems-problemphyxmini0100-prismvertex}
  \lean{PhyXMiniProblems.ProblemPhyXMini0100.PrismVertex}
  The three vertices identified by their angle labels in the figure
\end{definition}
\begin{proof}
  This is the declared model definition of Prism Vertex.
\end{proof}
\begin{definition}[Prism Face]
  \label{def:physics:phyx-mini-0100:phyxminiproblems-problemphyxmini0100-prismface}
  \lean{PhyXMiniProblems.ProblemPhyXMini0100.PrismFace}
  The three faces of the triangular prism cross-section
\end{definition}
\begin{proof}
  This is the declared model definition of Prism Face.
\end{proof}
\begin{definition}[Prism Liquid Setup]
  \label{def:physics:phyx-mini-0100:phyxminiproblems-problemphyxmini0100-prismliquidsetup}
  \lean{PhyXMiniProblems.ProblemPhyXMini0100.PrismLiquidSetup}
  \uses{def:physics:phyx-mini-0100:phyxminiproblems-problemphyxmini0100-prismvertex, def:physics:phyx-mini-0100:phyxminiproblems-problemphyxmini0100-prismface}
  The geometric and optical quantities belonging to the prism experiment. internalRayAngleToHypotenuseRadians is measured from the hypotenuse itself, whereas incidenceAtHypotenuseFromNormalRadians is measured from the normal to that face. No value for the requested liquid index is stored here
\end{definition}
\begin{proof}
  This is the declared model definition of Prism Liquid Setup.
\end{proof}
\begin{definition}[Matches Prism Figure]
  \label{def:physics:phyx-mini-0100:phyxminiproblems-problemphyxmini0100-matchesprismfigure}
  \lean{PhyXMiniProblems.ProblemPhyXMini0100.MatchesPrismFigure}
  \uses{def:physics:phyx-mini-0100:phyxminiproblems-problemphyxmini0100-prismliquidsetup}
  The labeled faces, vertex angles, and normal-entry datum read from the figure
\end{definition}
\begin{proof}
  This is the declared model definition of Matches Prism Figure.
\end{proof}
\begin{definition}[Matches Prism Index Readout]
  \label{def:physics:phyx-mini-0100:phyxminiproblems-problemphyxmini0100-matchesprismindexreadout}
  \lean{PhyXMiniProblems.ProblemPhyXMini0100.MatchesPrismIndexReadout}
  \uses{def:physics:phyx-mini-0100:phyxminiproblems-problemphyxmini0100-prismliquidsetup}
  The supplied dimensionless refractive-index readout n\_prism = 1.56
\end{definition}
\begin{proof}
  This is the declared model definition of Matches Prism Index Readout.
\end{proof}
\begin{definition}[Satisfies Prism Ray Geometry]
  \label{def:physics:phyx-mini-0100:phyxminiproblems-problemphyxmini0100-satisfiesprismraygeometry}
  \lean{PhyXMiniProblems.ProblemPhyXMini0100.SatisfiesPrismRayGeometry}
  \uses{def:physics:phyx-mini-0100:phyxminiproblems-problemphyxmini0100-prismliquidsetup}
  The elementary perpendicular-line geometry followed by the internal ray. Normal entry makes the ray-to-hypotenuse angle complementary to the 60° vertex angle. Incidence from the hypotenuse normal is in turn complementary to the ray-to-face angle. These are governing geometric relations, not a statement of the requested maximum refractive index
\end{definition}
\begin{proof}
  This is the declared model definition of Satisfies Prism Ray Geometry.
\end{proof}
\begin{definition}[Has Physical Prism Parameters]
  \label{def:physics:phyx-mini-0100:phyxminiproblems-problemphyxmini0100-hasphysicalprismparameters}
  \lean{PhyXMiniProblems.ProblemPhyXMini0100.HasPhysicalPrismParameters}
  \uses{def:physics:phyx-mini-0100:phyxminiproblems-problemphyxmini0100-prismliquidsetup}
  Positivity and principal acute branches for the physical optical data
\end{definition}
\begin{proof}
  This is the declared model definition of Has Physical Prism Parameters.
\end{proof}
\begin{definition}[Permits Total Internal Reflection]
  \label{def:physics:phyx-mini-0100:phyxminiproblems-problemphyxmini0100-permitstotalinternalreflection}
  \lean{PhyXMiniProblems.ProblemPhyXMini0100.PermitsTotalInternalReflection}
  \uses{def:physics:phyx-mini-0100:phyxminiproblems-problemphyxmini0100-prismliquidsetup}
  For a candidate liquid index n\_liquid, the limiting Snell-law condition is n\_liquid sin(π/2) ≤ n\_prism sin(θ\_incidence). The endpoint is included because the question asks for the largest permissible index: equality is the critical, tangential-transmission boundary convention used by the textbook calculation. The strict density ordering separately records that total internal reflection is from the prism toward the less optically dense liquid
\end{definition}
\begin{proof}
  This is the declared model definition of Permits Total Internal Reflection.
\end{proof}
\begin{definition}[admissible Liquid Refractive Indices]
  \label{def:physics:phyx-mini-0100:phyxminiproblems-problemphyxmini0100-admissibleliquidrefractiveindices}
  \lean{PhyXMiniProblems.ProblemPhyXMini0100.admissibleLiquidRefractiveIndices}
  \uses{def:physics:phyx-mini-0100:phyxminiproblems-problemphyxmini0100-prismliquidsetup, def:physics:phyx-mini-0100:phyxminiproblems-problemphyxmini0100-permitstotalinternalreflection}
  All positive liquid-index readouts admitted by the limiting TIR criterion
\end{definition}
\begin{proof}
  This is the declared model definition of admissible Liquid Refractive Indices.
\end{proof}
\begin{definition}[critical Liquid Index Candidate]
  \label{def:physics:phyx-mini-0100:phyxminiproblems-problemphyxmini0100-criticalliquidindexcandidate}
  \lean{PhyXMiniProblems.ProblemPhyXMini0100.criticalLiquidIndexCandidate}
  \uses{def:physics:phyx-mini-0100:phyxminiproblems-problemphyxmini0100-prismliquidsetup}
  The liquid index obtained by putting the transmitted critical ray tangent to the interface in Snell's law. This names a candidate boundary value; it does not assert that the candidate is admissible or greatest
\end{definition}
\begin{proof}
  This is the declared model definition of critical Liquid Index Candidate.
\end{proof}
\begin{definition}[Answer Choice]
  \label{def:physics:phyx-mini-0100:phyxminiproblems-problemphyxmini0100-answerchoice}
  \lean{PhyXMiniProblems.ProblemPhyXMini0100.AnswerChoice}
  The labels printed beside the four multiple-choice values
\end{definition}
\begin{proof}
  This is the declared model definition of Answer Choice.
\end{proof}
\begin{definition}[answer Choice Refractive Index]
  \label{def:physics:phyx-mini-0100:phyxminiproblems-problemphyxmini0100-answerchoicerefractiveindex}
  \lean{PhyXMiniProblems.ProblemPhyXMini0100.answerChoiceRefractiveIndex}
  \uses{def:physics:phyx-mini-0100:phyxminiproblems-problemphyxmini0100-answerchoice}
  The dimensionless liquid-index value printed for each answer choice
\end{definition}
\begin{proof}
  This is the declared model definition of answer Choice Refractive Index.
\end{proof}
\begin{definition}[Is Closest Answer Choice]
  \label{def:physics:phyx-mini-0100:phyxminiproblems-problemphyxmini0100-isclosestanswerchoice}
  \lean{PhyXMiniProblems.ProblemPhyXMini0100.IsClosestAnswerChoice}
  \uses{def:physics:phyx-mini-0100:phyxminiproblems-problemphyxmini0100-answerchoice, def:physics:phyx-mini-0100:phyxminiproblems-problemphyxmini0100-answerchoicerefractiveindex}
  A displayed choice is uniquely closest to an exact liquid-index value
\end{definition}
\begin{proof}
  This is the declared model definition of Is Closest Answer Choice.
\end{proof}
\begin{lemma}[hypotenuse incidence is sixty degrees]
  \label{lem:physics:phyx-mini-0100:phyxminiproblems-problemphyxmini0100-hypotenuse-incidence-is-sixty-degrees}
  \lean{PhyXMiniProblems.ProblemPhyXMini0100.hypotenuse_incidence_is_sixty_degrees}
  \uses{def:physics:phyx-mini-0100:phyxminiproblems-problemphyxmini0100-prismliquidsetup, def:physics:phyx-mini-0100:phyxminiproblems-problemphyxmini0100-matchesprismfigure, def:physics:phyx-mini-0100:phyxminiproblems-problemphyxmini0100-satisfiesprismraygeometry}
  Normal entry through the short leg of the depicted prism makes the incidence angle at the hypotenuse equal to 60°
\end{lemma}
\begin{proof}
  The conclusion follows from the governing relations and figure readouts named in the statement of hypotenuse incidence is sixty degrees.
\end{proof}
\begin{lemma}[critical liquid index exact]
  \label{lem:physics:phyx-mini-0100:phyxminiproblems-problemphyxmini0100-critical-liquid-index-exact}
  \lean{PhyXMiniProblems.ProblemPhyXMini0100.critical_liquid_index_exact}
  \uses{def:physics:phyx-mini-0100:phyxminiproblems-problemphyxmini0100-prismliquidsetup, def:physics:phyx-mini-0100:phyxminiproblems-problemphyxmini0100-matchesprismfigure, def:physics:phyx-mini-0100:phyxminiproblems-problemphyxmini0100-matchesprismindexreadout, def:physics:phyx-mini-0100:phyxminiproblems-problemphyxmini0100-satisfiesprismraygeometry, def:physics:phyx-mini-0100:phyxminiproblems-problemphyxmini0100-criticalliquidindexcandidate}
  With n\_prism = 1.56 and θ\_incidence = 60°, the critical Snell candidate is 1.56 sin(60°) = 39 √3 / 50
\end{lemma}
\begin{proof}
  The conclusion follows from the governing relations and figure readouts named in the statement of critical liquid index exact.
\end{proof}
% --- Archon declaration topology end ---
% --- Archon physics formalization source end ---
